\documentclass[12pt, twoside]{article}
%\documentclass[12pt, twoside]{article}
\usepackage[letterpaper, margin=1in, headsep=0.2in]{geometry}
\setlength{\headheight}{0.6in}
%\usepackage[english]{babel}
\usepackage[utf8]{inputenc}
\usepackage{microtype}
\usepackage{amsmath}
\usepackage{amssymb}
%\usepackage{amsfonts}
\usepackage[nomessages]{fp} %\FPeval{\var-name}{2*sin(pi/6)}
\usepackage{siunitx} %units in math. eg 20\milli\meter
\usepackage{yhmath} % for arcs, overparenth command
\usepackage{tikz} %graphics
\usetikzlibrary{quotes, angles, arrows, arrows.meta}
\usepackage{pgfplots}
\usepgfplotslibrary{fillbetween}
\pgfplotsset{compat=1.16}
\usepackage{graphicx} %consider setting \graphicspath{{images/}}
\usepackage{parskip} %no paragraph indent
\usepackage{enumitem}
\usepackage{multicol}
\usepackage{venndiagram}

\usepackage{fancyhdr}
\pagestyle{fancy}
\fancyhf{}
\renewcommand{\headrulewidth}{0pt} % disable the underline of the header
\raggedbottom
\hfuzz=2mm %suppresses overfull box warnings

\usepackage{hyperref}

\fancyhead[LE]{\thepage}
\fancyhead[RO]{Solutions}
\fancyhead[LO]{La Scuola d'Italia / Dr.\ Huson / IB Math 11 \\* 5 October 2026}

\begin{document}
\subsubsection*{1.9 Pre-Quiz: Arithmetic and Geometric Sequences --- Solutions}

\emph{Review set, not scored: worked solutions only, no marks. Calculator
allowed. Answers are exact or correct to three significant figures.}

\begin{enumerate}[itemsep=0.3cm]

%--- Problem 1: Arithmetic sequence; d, u_n, u_40, which values are terms ---
\item[1.]

$12,\; 19,\; 26,\; 33,\; \ldots$

\medskip

\textbf{(a)} $d = 19 - 12 = 7$

\medskip

\textbf{(b)} $u_{n} = u_{1} + (n - 1)d$

$u_{n} = 12 + 7(n - 1)$

\emph{Accept $u_{n} = 7n + 5$.}

\medskip

\textbf{(c)} $u_{40} = 12 + 7(40 - 1) = 12 + 273 = 285$

\medskip

\textbf{(d)} $12 + 7(n - 1) = 250 \;\Rightarrow\; 7(n - 1) = 238
\;\Rightarrow\; n - 1 = 34 \;\Rightarrow\; n = 35$

$n$ is a whole number, so $250$ is the $35^{\text{th}}$ term.

\medskip

\textbf{(e)} $12 + 7(n - 1) = 300 \;\Rightarrow\; 7(n - 1) = 288
\;\Rightarrow\; n - 1 = 41.14\ldots \;\Rightarrow\; n = 42.14\ldots$

$n$ is not a whole number, so $300$ is not a term.

\emph{Equivalently: $u_{42} = 299$ and $u_{43} = 306$, and $300$ lies
between them.}

\bigskip \hrule \bigskip

%--- Problem 2: Two given terms -> u_1 and d; first negative term; terms in k ---
\item[2.]

$u_{5} = 31$, \ $u_{12} = 3$.

\medskip

\textbf{(a)} From $u_{5}$ to $u_{12}$ is $7$ steps of $d$:
$7d = 3 - 31 = -28 \;\Rightarrow\; d = -4$

$u_{5} = u_{1} + 4d \;\Rightarrow\; 31 = u_{1} - 16 \;\Rightarrow\;
u_{1} = 47$

\emph{Or solve the simultaneous equations $u_{1} + 4d = 31$,
$u_{1} + 11d = 3$ on the GDC.}

\medskip

\textbf{(b)} $u_{n} = 47 - 4(n - 1)$

\emph{Accept $u_{n} = 51 - 4n$.}

\medskip

\textbf{(c)} $47 - 4(n - 1) < 0 \;\Rightarrow\; 4(n - 1) > 47
\;\Rightarrow\; n - 1 > 11.75 \;\Rightarrow\; n > 12.75$, so $n = 13$.

$u_{12} = 3$ and $u_{13} = 47 - 4(13 - 1) = -1$. The first negative term is $-1$ (the
$13^{\text{th}}$ term).

\emph{The question asks for the term, not $n$: the answer is $-1$.}

\medskip

\textbf{(d)} Equal differences: \ $3k - (k + 2) = (4k + 1) - 3k$

$2k - 2 = k + 1 \;\Rightarrow\; k = 3$

First three terms: $5,\; 9,\; 13$ \quad ($d = 4$)

\medskip

\textbf{(e)} $u_{20} = 5 + 4(20 - 1) = 5 + 76 = 81$

\newpage

%--- Problem 3: Arithmetic, geometric or neither; geometric u_n, u_10, first term over 1000; two given terms ---
\item[3.]

\textbf{(a)}

\begin{tabular}{rll}
(i) & $3,\; 7,\; 11,\; 15$ & arithmetic, $d = 4$ \\
(ii) & $3,\; 6,\; 12,\; 24$ & geometric, $r = 2$ \\
(iii) & $3,\; 4,\; 6,\; 9$ & neither (differences $1, 2, 3$; ratios
$1.33, 1.5, 1.5$) \\
(iv) & $80,\; 20,\; 5,\; 1.25$ & geometric, $r = \frac{1}{4} = 0.25$ \\
\end{tabular}

\emph{Test arithmetic by subtracting consecutive terms, geometric by
dividing. In (iii) neither the differences nor the ratios are all
equal.}

\medskip

\textbf{(b)} $r = \dfrac{9}{6} = 1.5$ \quad \emph{(check:
$\dfrac{13.5}{9} = 1.5$)}

\medskip

\textbf{(c)} $u_{n} = u_{1}\,r^{\,n - 1} = 6(1.5)^{n - 1}$

\medskip

\textbf{(d)} $u_{10} = 6(1.5)^{9} = 230.66\ldots \approx 231$

\emph{The exponent is $9$, not $10$. $6(1.5)^{10} = 346$ is the
$11^{\text{th}}$ term.}

\medskip

\textbf{(e)} $6(1.5)^{n - 1} > 1000$

$u_{13} = 778.5$ \ ($< 1000$); \quad $u_{14} = 1167.7$ \ ($> 1000$)

$n = 14$

\emph{Or solve $1.5^{\,n - 1} > 166.7$ on the GDC: $n - 1 > 12.6$, so
$n = 14$. The answer is a whole number.}

\medskip

\textbf{(f)} From $u_{2}$ to $u_{5}$ is $3$ steps of $\times r$:
$r^{3} = \dfrac{324}{12} = 27 \;\Rightarrow\; r = 3$

$u_{1} = \dfrac{u_{2}}{r} = \dfrac{12}{3} = 4$

\emph{Check: $4,\; 12,\; 36,\; 108,\; 324$.}

\newpage

%--- Problem 4: Applications -- arithmetic vs geometric salary; compound interest and real value ---
\item[4.]

Job A: arithmetic, $u_{1} = 42\,000$, $d = 1500$. \ Job B: geometric,
$u_{1} = 40\,000$, $r = 1.05$.

\medskip

\textbf{(a)} Job A: $42\,000 + 2 \times 1500 = $ \$$45\,000$

Job B: $40\,000 \times 1.05^{2} = $ \$$44\,100$

\medskip

\textbf{(b)} Job A: $A_{n} = 42\,000 + 1500(n - 1) = 40\,500 + 1500n$

Job B: $B_{n} = 40\,000(1.05)^{n - 1}$

\medskip

\textbf{(c)} Compare the two salaries year by year:

\begin{tabular}{ccc}
Year & Job A & Job B \\
\hline
$4$ & \$$46\,500$ & \$$46\,305$ \\
$5$ & \$$48\,000$ & \$$48\,620.25$ \\
\end{tabular}

Year $5$

\emph{Use a GDC table of both expressions. Show the two adjacent years:
year $4$ is what shows that year $5$ is the \emph{first}.}

\medskip

\textbf{(d)} $3000 \times 1.03^{8} = $ \$$3800.31$ \quad (\$$3800$ to
3 s.f.)

\emph{TVM solver: $N = 8$, $I\% = 3$, $PV = -3000$, $PMT = 0$,
$P/Y = C/Y = 1$.}

\medskip

\textbf{(e)} Real value $= \dfrac{3800.31}{1.02^{8}} = $ \$$3243.53$
\quad (\$$3240$ to 3 s.f.)

\emph{The real-interest-rate method is also accepted: $3 - 2 = 1\%$ per
year, $3000 \times 1.01^{8} = $ \$$3248.57$ (\$$3250$ to 3 s.f.).}

\emph{The real value must be less than the answer to (d). Multiplying
by $1.02^{8}$ is the most common error.}

\end{enumerate}
\end{document}
